\chapter{الجمع والطرح}\label{ch:g3:addsub}

الجمع ضمٌّ، والطرح أخذٌ أو بحث عمّا ينقص. يدرّب هذا الفصل على
طريقتَي العمود --- بما فيهما من احتفاظ واستلاف مشهورين --- وعلى
الحيل الذهنية التي تغني عن الأعمدة في كثير من
الأحيان.

\section{الجمع بالأعمدة}

\begin{method}[الجمع في أعمدة]\label{met:g3:addsub:add}
\begin{enumerate}
  \item اُكتب العددين أحدهما تحت الآخر، الآحاد تحت الآحاد،
        والعشرات تحت العشرات (والمحاذاة من \emph{اليمين})؛
  \item اِجمع الآحاد؛ فإن تجاوز \omterm{def:g1:addition:def}{المجموع} $9$، اُكتب رقم آحاده
        و\emph{احتفظ} بالعشرة للعمود التالي (رقم $1$ صغير)؛
  \item وواصل مع العشرات، ثم المئات، \dots، مضيفًا المحتفظ به
        في كل مرة.
\end{enumerate}
\end{method}

\begin{example}\label{ex:g3:addsub:add}
احسب $457 + 368$:
\[
\begin{array}{r}
{\scriptstyle 1}\ \ {\scriptstyle 1}\ \ \phantom{0} \\[-3pt]
4\,5\,7 \\
+\ 3\,6\,8 \\
\hline
8\,2\,5
\end{array}
\]
الآحاد: $7 + 8 = 15$: اُكتب $5$ واحتفظ برقم $1$. والعشرات:
$5 + 6 + 1 = 12$: اُكتب $2$ واحتفظ برقم $1$. والمئات: $4 + 3 + 1 = 8$.
والنتيجة: $825$.
\end{example}

\begin{example}[الجمع الذهني]\label{ex:g3:addsub:mental}
اِبحث عن الأعداد الودودة قبل أن تلجأ إلى الأعمدة:
\begin{itemize}
  \item \emph{اِصنع عشرة}: $47 + 8 = 47 + 3 + 5 = 50 + 5 = 55$؛
  \item \emph{اقفز بالأجزاء}: $256 + 130 = 256 + 100 + 30 = 386$؛
  \item \emph{أعد الترتيب}: $17 + 59 + 3 = (17 + 3) + 59 = 79$.
\end{itemize}
\end{example}

\begin{omfigure}
\begin{tikzpicture}[scale=1.05]
  \draw[-Stealth] (-0.2,0) -- (10.4,0);
  \foreach \x/\l in {0.56/{256}, 5.56/{356}, 7.06/{386}}
    { \draw (\x,-0.08) -- (\x,0.08);
      \node[below, font=\small] at (\x,-0.1) {$\l$}; }
  \draw[-Stealth, omDef, thick] (0.56,0.25) .. controls (2.8,1.15) ..
    (5.56,0.25);
  \node[omDef, font=\small] at (3.05,1.05) {$+100$};
  \draw[-Stealth, omThm, thick] (5.56,0.25) .. controls (6.3,0.8) ..
    (7.06,0.25);
  \node[omThm, font=\small] at (6.3,0.85) {$+30$};
\end{tikzpicture}

{\small الجمع الذهني قفزات على مستقيم الأعداد: $256 + 130$ في
قفزتين، المئات أولًا ثم العشرات.}
\end{omfigure}

\section{الطرح بالأعمدة}

\begin{method}[الطرح في أعمدة]\label{met:g3:addsub:sub}
\begin{enumerate}
  \item اُكتب العدد الأكبر في الأعلى، محاذيًا من اليمين؛
  \item واطرح عمودًا عمودًا، بدءًا من الآحاد؛
  \item فإن كان الرقم الأعلى أصغر من أن يكفي، \emph{فاستلف}: خذ
        واحدًا من المرتبة التالية في السطر الأعلى (فيساوي عشرة في
        العمود الحالي)؛
  \item \emph{تحقّق}: النتيجة زائد المطروح يجب أن تعيد
        العدد الأعلى.
\end{enumerate}
\end{method}

\begin{example}\label{ex:g3:addsub:sub}
احسب $603 - 147$:
\[
\begin{array}{r}
6\,0\,3 \\
-\ 1\,4\,7 \\
\hline
4\,5\,6
\end{array}
\]
الآحاد: $3 - 7$ غير ممكن؛ فاستلف عبر المراتب: فيصير $603$
هو $5$ مئات و $9$ عشرات و $13$ واحدًا. ثم $13 - 7 = 6$؛
والعشرات $9 - 4 = 5$؛ والمئات $5 - 1 = 4$. والنتيجة: $456$.
والتحقّق: $456 + 147 = 603$.
\end{example}

\begin{omfigure}
\begin{tikzpicture}[scale=0.82]
  % الحالة الأولى: 6 مئات و0 عشرات و3 آحاد
  \node[font=\small] at (2.5,2.25) {$603$ قبل الاستلاف};
  \foreach \i in {0,...,5}
    \draw[omProp, fill=omProp!25] (0.62*\i,0.7) rectangle
      (0.62*\i+0.44,1.7);
  \node[below, font=\footnotesize] at (1.7,0.55) {$6$ مئات};
  \node[font=\footnotesize] at (4.35,1.2) {$0$ عشرات};
  \foreach \i in {0,1,2} \fill[omDef] (5.4+0.4*\i,1.2) circle (0.13);
  \node[below, font=\footnotesize] at (5.8,0.55) {$3$ آحاد};
  % الحالة الثانية: 5 مئات و9 عشرات و13 واحدًا
  \begin{scope}[yshift=-3.3cm]
  \node[font=\small] at (2.5,2.25) {$603$ بعد الاستلاف};
  \foreach \i in {0,...,4}
    \draw[omProp, fill=omProp!25] (0.62*\i,0.7) rectangle
      (0.62*\i+0.44,1.7);
  \node[below, font=\footnotesize] at (1.4,0.55) {$5$ مئات};
  \foreach \i in {0,...,8}
    \draw[omThm, very thick] (3.5+0.22*\i,0.7) -- (3.5+0.22*\i,1.7);
  \node[below, font=\footnotesize] at (4.4,0.55) {$9$ عشرات};
  \foreach \i in {0,...,6} \fill[omDef] (6.1+0.4*\i,1.45) circle (0.13);
  \foreach \i in {0,...,5} \fill[omDef] (6.1+0.4*\i,0.95) circle (0.13);
  \node[below, font=\footnotesize] at (7.3,0.55) {$13$ واحدًا};
  \end{scope}
\end{tikzpicture}

{\small استلاف $603 - 147$ بالصورة: مئة واحدة تُفكّ إلى عشر
عشرات، وواحدة من تلك العشرات إلى عشرة آحاد. المقدار نفسه في ثوب
جديد --- فيصير $6\,\vert\,0\,\vert\,3$ هو
$5\,\vert\,9\,\vert\,13$، ويصبح الطرح ممكنًا في كل عمود.}
\end{omfigure}

\begin{example}[الطرح بحثًا عن الجزء الناقص]\label{ex:g3:addsub:missing}
«عند ليلى $35$ كرة، وعند عمر $52$: بكم يزيد عمر؟» هو
الطرح $52 - 35$. وذهنيًّا، عُدّ \emph{إلى الأمام} من $35$:
من $35 \to 40$ مقداره $5$، ثم من $40 \to 52$ مقداره $12$؛ وفي المجموع
$5 + 12 = 17$. والعدّ إلى الأمام أسهل من الاستلاف في الغالب.
\end{example}

\begin{omfigure}
\begin{tikzpicture}[scale=1.05]
  \draw[-Stealth] (-0.2,0) -- (10.4,0);
  \foreach \x/\l in {1.0/{35}, 3.5/{40}, 9.5/{52}}
    { \draw (\x,-0.08) -- (\x,0.08);
      \node[below, font=\small] at (\x,-0.1) {$\l$}; }
  \draw[-Stealth, omDef, thick] (1.0,0.25) .. controls (2.2,0.9) ..
    (3.5,0.25);
  \node[omDef, font=\small] at (2.25,0.92) {$+5$};
  \draw[-Stealth, omThm, thick] (3.5,0.25) .. controls (6.5,1.2) ..
    (9.5,0.25);
  \node[omThm, font=\small] at (6.5,1.1) {$+12$};
\end{tikzpicture}

{\small حساب $52 - 35$ بالعدّ إلى الأمام: إلى $40$ الودود
أولًا، ثم إلى $52$. والجواب هو \omterm{def:g1:addition:def}{مجموع} القفزتين:
$5 + 12 = 17$.}
\end{omfigure}

\section{رتبة القدر}

\begin{method}[التقدير قبل الحساب]\label{met:g3:addsub:estimate}
قبل أي حساب بالأعمدة، قرِّب الأعداد إلى أعداد ودودة واحسب
\emph{تقديرًا}. وبعد الحساب، قارن: فنتيجة بعيدة عن التقدير تنبّه
إلى خطأ.
\end{method}

\begin{example}
في $487 + 312$: التقدير هو $500 + 300 = 800$. ونتيجة الأعمدة،
$799$، قريبة منه: فهي معقولة. ولو وجدت $1\,099$
(بسبب محاذاة خاطئة)، لكان التقدير قد نبّهك.
\end{example}

\section{تمارين}

\begin{exercise}[$\star$]\label{exo:g3:addsub:1}
احسب بالأعمدة: $346 + 253$؛ \quad $478 + 265$؛ \quad
$1\,536 + 2\,876$.
\end{exercise}

\begin{exercise}[$\star$]\label{exo:g3:addsub:2}
احسب بالأعمدة، ثم تحقّق بعملية جمع: $587 - 234$؛ \quad
$805 - 348$؛ \quad $1\,000 - 463$.
\end{exercise}

\begin{exercise}[$\star$]\label{exo:g3:addsub:3}
احسب ذهنيًّا بصنع العشرات: $38 + 7$؛ \quad $56 + 9$؛ \quad
$45 + 28$ (أضف $30$، ثم اطرح $2$).
\end{exercise}

\begin{exercise}[$\star$]\label{exo:g3:addsub:4}
احسب ذهنيًّا بالقفزات: $340 + 220$؛ \quad $675 + 210$؛ \quad
$512 - 300$؛ \quad $840 - 220$.
\end{exercise}

\begin{exercise}[$\star$]\label{exo:g3:addsub:5}
احسب $63 - 47$ بالعدّ إلى الأمام (كما في
\cref{ex:g3:addsub:missing})، ثم تحقّق بالأعمدة.
\end{exercise}

\begin{exercise}[$\star$]\label{exo:g3:addsub:6}
قدِّر أولًا (بالتقريب إلى المئات)، ثم احسب بالضبط:
$492 + 218$؛ \quad $913 - 388$.
\end{exercise}

\begin{exercise}[$\star$]\label{exo:g3:addsub:7}
في مدرسة $348$ فتاة و $296$ فتى. فكم تلميذًا فيها كلها؟
اُكتب العملية، واحسب، وأجب بجملة كاملة.
\end{exercise}

\begin{exercise}[$\star$]\label{exo:g3:addsub:8}
في كتاب $224$ صفحة، قرأت نورة منها $178$. فكم صفحة بقيت لها
لتقرأها؟
\end{exercise}

\begin{exercise}[$\star$]\label{exo:g3:addsub:9}
أكمل الفراغات (اعمل عمودًا عمودًا):
\[
\begin{array}{r}
3\,?\,7 \\
+\ ?\,2\,? \\
\hline
8\,6\,2
\end{array}
\]
فهل هناك طريقة واحدة فقط لملء الفراغات الثلاثة؟ اِشرح.
\end{exercise}

\begin{exercise}[$\star\star$]\label{exo:g3:addsub:10}
باع مخبز $145$ كرواسانًا صباحًا و $89$ بعد الظهر؛ وبقي $17$
كرواسانًا لم يُبع عند الإغلاق. فكم كرواسانًا خُبز في ذلك اليوم؟
(خطوتان.)
\end{exercise}

\begin{exercise}[$\star\star$]\label{exo:g3:addsub:11}
يقول طارق: «لحساب $500 - 199$، أحسب $500 - 200$ ثم
أضيف $1$». أهو محقّ؟ احسب بالطريقتين، واشرح حيلته.
\end{exercise}
